\section{Diagnostic Results Draft}

The current experiments should be interpreted as diagnostics for an inspectable research scaffold.
They test whether the implementation can expose differences between simple heuristics, model-free learning, and the decomposed strategy loop.
They do not establish robust multi-agent reasoning or equilibrium convergence.

\subsection{Gridworld Baselines}

The May 15 public snapshot compares four baselines on the default 10 by 10 attacker-defender gridworld.
The direct-to-goal heuristic reaches the goal in every evaluation episode, with win rate $1.000$ and mean episode return $0.910$.
The PPO-lite baseline, random policy, and early Day-2 strategy loop all have win rate $0.000$ under the short diagnostic budget.
Their mean returns are $-1.180$, $-1.208$, and $-1.220$, respectively.

This result is useful because it falsifies an easy overclaim: the current latent strategy loop is not yet a strong controller for the default gridworld.
For a methods paper, the baseline table should be framed as evidence that the harness can detect trivial but important baselines, not as evidence that the energy-based method is already competitive.

\subsection{Payoff and Strategy Diagnostics}

The named-strategy payoff matrix provides a compact diagnostic for how heuristic strategy embeddings behave against sampled opponent responses.
It should be plotted as a heatmap and described as sampled-response evaluation.
The matrix is not an equilibrium computation, and the worst sampled response is not an exact best response over the full policy space.

\subsection{Ablation Protocol}

The paper-facing ablation suite compares Langevin EBM sampling against Gaussian latent sampling, buffer positives against random positives, robustness-aware selection against average-value selection, different candidate counts, different opponent sample counts, and world-model fitting on or off.
The purpose of these ablations is to identify which parts of the decomposition change diagnostic metrics such as return, win rate, buffer diversity, robustness, and the exploitability proxy.
The purpose is not to claim state-of-the-art performance.

\section{Limitations Draft}

All game-theoretic metrics in the current scaffold are approximate.
In particular, the reported \texttt{exploitability\_proxy} is computed against a finite sampled set of heuristic opponent responses.
It is not Nash exploitability, it is not a CFR or PSRO estimate, and it should not be used as evidence of equilibrium quality.

The energy model is currently a small MLP trained with a contrastive divergence-style objective over buffer positives and sampled negatives.
It does not use persistent contrastive divergence, a learned proposal distribution, or score matching in the active training loop.
Reported strategy diversity can be low, and early public snapshots include cases where diversity collapses.

The world model is fit to one-step transitions but is not yet used for imagined rollouts or model-based strategy evaluation.
Therefore, any claim about model-based reasoning should be deferred until the rollout interface is implemented and evaluated.

The PPO baselines are intentionally lightweight.
They are useful as live learning baselines and sanity checks, but stronger PPO budgets and multi-seed learning curves are needed before comparing learning efficiency.

The multi-evader pursuit infrastructure is valuable for trace inspection and empirical-game diagnostics, but it is currently a supplementary track.
It should not be presented as a second EBM-driven domain unless latent strategy generation and sampled-response selection are explicitly integrated into that environment.
